\chapter{Ordered Sets}
\label{chap:ordered-sets}

% Closed question: What is the smallest admissible mark a ledger can carry?

\begin{chapteressay}

Every measurement assigns a number by counting. The yardstick says how many
ticks fit across the room. The clock says how many tocks have passed.
Distance is a count of length ticks; duration is a count of time ticks.
Speed is the ratio of two counts. The instrument is the device that
performs the counting; the number is the count it produces; the
measurement is the number interpreted as the count under a stated
calibration.

This is the ground example of the mathematical gauge. Everything that
follows --- partial orders, lattices, splines, transports, curvatures,
invariants, closures --- emerges from the structure of counting and the
ratios of counts. The number is not the world. The number is the count of
ticks; the world is whatever the ticks were placed to detect.

A counter operates by distinguishing one tick from the next. The two ticks
must be \emph{distinct} in some way the counter can recognize. If two ticks
were identical to the counter, they would be one tick, not two. The
counting requires distinction. The distinction is the first mathematical
event in any measurement.

Before the count, there is the act of distinguishing. Before the act of
distinguishing, there is the set whose elements are being distinguished.
The smallest unit of mathematical structure required for measurement is a
set with a relation that says which pairs of elements are admissibly
distinguishable. The relation must be reflexive (every element is
indistinguishable from itself), antisymmetric (if $a$ is indistinguishable
from $b$ and $b$ is indistinguishable from $a$, then $a$ and $b$ are the
same element), and transitive (if $a$ is indistinguishable from $b$ and
$b$ is indistinguishable from $c$, then $a$ is indistinguishable from
$c$).

These three conditions define a partial equivalence. Adding decidability
(the relation can be checked algorithmically) gives an admissible
distinction relation. The chapter's claim is that this minimum is the
foundational structure on which all subsequent mathematics rests. Without
the distinction relation, there is no counting. Without counting, there
is no measurement. Without measurement, there are no numbers.

This is not a reduction of mathematics to bookkeeping. The chapter's
argument is the reverse: mathematics, when audited at its foundation,
already contains the structure of counting and distinction. The
sophisticated machinery of analysis (Cauchy completion, limits, residues)
emerges from the foundational structure, not as a separate invention.
Every classical theorem can be traced back to the distinction relation
on some underlying set.

Take a Roman numeral. The mark VII carries the count seven. A
fifth-marker V, then two strokes I. The V is distinct from the
strokes; the strokes are distinct from one another. The count is
the value of the V plus the number of strokes, totaled in the
conventions of Roman arithmetic. The same count is also carried by
7 (Arabic), by 111 (binary), by \emph{sieben} (German). Each
notation is a different carrier; each carrier denotes the same
count.

The carrier is the symbol; the count is the value. They are not the
same. A mark in the ledger is a pair: carrier and value. Replacing the
carrier with a different notation does not change the value. The
mathematical content of the mark is the value; the operational form is
the carrier.

The distinction relation operates on the values. Two values are
distinguishable when they are different counts. The carrier may differ
without the values differing (different notations of the same count),
or the values may differ without the carriers obviously differing
(easily confused notations of different counts). The adequacy of a
carrier system is the requirement that distinct values receive distinct
carriers, and that the distinction can be decided algorithmically.

The smallest admissible mark a ledger can carry is therefore a triple:
a set element (the value), a carrier (the symbol that names it), and a
distinction relation (which admits the equality check that distinguishes
this carrier from other carriers in the system). The triple is not
contingent: it is the minimum mathematical structure required for the
ledger to be a ledger.

The chapter develops the triple's consequences in five sections.

Section 1 introduces sets and relations as the foundational pair. The
example is RNA base-pair nesting, which generates a partial order on
positions from molecular geometry. The order is forced by the
chemistry, not imposed. The non-crossing condition on pairs is the
chapter's first instance of an admissible partial order.

Section 2 introduces the carrier. The example is the Roman numeral
VII as the carrier of the count seven. The distinction between
carrier and value is the chapter's first structural separation. The
DISTINGUISHABLE typeclass formalizes the adequacy of carriers.

Section 3 introduces partial order and distinguishability. The
example is the divisibility relation on the positive integers: a
partial order that is not total, with meet and join given by gcd and
lcm. The example illustrates that orders can capture structure that
total orders would collapse.

Section 4 introduces Cauchy sequences and limits. The example is the
decimal expansion of $\pi$ as a Cauchy sequence of rationals whose
limit is irrational. The Cauchy completion of the rationals is the
mathematical structure forced by admitting limits.

Section 5 introduces the residue. The example is polynomial long
division: dividing $x^3 - 2$ by $x - 1$ leaves remainder $-1$. The
residue is part of the record; it is not an error or a remainder to
be discarded.

The second concrete example for the chapter is the archive accession
desk. The archivist gives every incoming document a unique accession
number. The number is the carrier; the document is the value. The
distinction relation is the database's uniqueness check. Multiple
partial orders are available (search by author, by date, by topic).
Supplementary additions to a collection are Cauchy refinements.
Unresolved metadata are residues that the catalog carries
explicitly. The accession desk is the chapter's structure
institutionalized.

The chapter's closed question is:

\begin{center}
\emph{What is the smallest admissible mark a ledger can carry?}
\end{center}

The chapter's answer is: a triple of value, carrier, and partial
order, with the adequacy condition that the carrier admits a
decision procedure for equality. Anything less is not yet a
ledger entry. The triple is the foundation on which subsequent
chapters build.

\end{chapteressay}

\section{Sets and Relations}
\label{se:sets-and-relations}

A set is a collection of distinguishable objects. The objects are the elements.
The collection itself is the set. Two sets are equal when they have the same
elements. A relation on a set is a subset of the Cartesian product of the set
with itself: the relation holds between two elements when their pair is in the
subset.

These two definitions, taken together, are minimal. They commit to almost
nothing: no arithmetic, no measurement, no geometry. They commit only that
elements can be distinguished from one another and that some pairs of elements
stand in a named relationship while others do not.

The smallest example is a two-element set with the equality relation: the only
related pairs are $(a,a)$ and $(b,b)$. The set $\{a, b\}$ has a relation
$\{(a,a), (b,b)\}$. The relation says: $a$ relates to itself, $b$ relates to
itself, $a$ and $b$ do not relate. This is enough to describe a record with
two marks.

\begin{phenom}{RNA Base-Pair Nesting}
\label{ph:rna-base-pair-nesting}

\PhStatement The hydrogen-bond pairings between bases in a folded RNA strand
form a partial order on the positions: if positions $i$ and $j$ are paired
and positions $k$ and $\ell$ are paired with $i < k < j$, then necessarily
$i < k < \ell < j$. The pairings nest; they do not cross.

\PhOrigin The non-crossing condition is a consequence of the molecular
geometry of RNA secondary structure. Crossing pairings would require the
backbone to fold through itself, which is sterically impossible for
typical secondary structures.

\PhObservation In a 76-nucleotide tRNA molecule, the base pairs from the
acceptor stem, the D arm, the anticodon arm, and the T arm all nest within
one another, producing the canonical clover-leaf structure.

\PhConstraint The pairing relation is a partial order on positions: the
pair $(i, j)$ is admissible only when no existing pair $(k, \ell)$ has
exactly one of $k, \ell$ between $i$ and $j$. Pairs are either nested
inside each other or disjoint.

\PhConsequence The set of pairings on a fixed RNA backbone forms a forest
under the nesting relation. The mathematics of secondary structure is
the combinatorics of non-crossing partitions.

\end{phenom}

The non-crossing pairings illustrate the chapter's first mathematical claim:
order can be forced by structure rather than imposed by convention. The
backbone's geometry does not say which bases must pair, but it does say that
whichever pairs occur must respect the non-crossing condition. The condition
is the partial order.

The smallest admissible mark a ledger can carry is, accordingly, not a
number. It is a pair: an element of a set and a relation that names its
position relative to the other elements. The element is the value; the
position is the order. Without the position, the value is unlocatable;
without the value, the position is empty.

A relation that is reflexive ($a$ relates to $a$ for every $a$), antisymmetric
(if $a$ relates to $b$ and $b$ relates to $a$, then $a = b$), and transitive
(if $a$ relates to $b$ and $b$ relates to $c$, then $a$ relates to $c$) is a
partial order. The three conditions are the minimum a relation must satisfy
to deserve the name. They are not optional. A relation lacking any one of
them is something weaker than an order.

In Lean, the partial order on natural numbers is the typeclass instance
\texttt{Number.le} in \texttt{Episode1.lean}. The instance provides the
function \texttt{le : Number $\to$ Number $\to$ Prop} together with proofs
of reflexivity, antisymmetry, and transitivity. The compiler checks each
proof at elaboration time. A proposed instance that fails any condition is
rejected.

The chapter's claim is that this triple --- set, relation, partial order ---
is the foundational structure on which all subsequent chapters build. Every
later development assumes it.

\section{The Carrier}
\label{se:the-carrier}

The Roman numeral VII carries the count seven. A fifth-mark V
followed by two parallel strokes I --- five plus one plus one ---
denotes the seventh natural number. The symbol is the carrier; the
count is the value. The symbol is the way the value enters the
ledger; the value is what the symbol designates.

The carrier is replaceable. The same count is also carried by the Arabic
numeral 7, by the binary string 111, by the German word \emph{sieben}, by
the prime factorization $7 = 7$, and by the seventh element of any
enumerated set. Each of these is a different symbol; each of these denotes
the same count. The count is invariant under the change of carrier; the
carrier is variable in the way one carries the count.

This distinction --- between the carrier and what it carries --- is the
chapter's second mathematical structure. Without the distinction, the
ledger would be a collection of symbols indistinguishable from a
collection of counts. With it, the ledger has a clear separation: marks
in the ledger are symbols; the symbols denote values; comparisons among
values do not depend on which symbols were used.

The adequacy of a carrier is the requirement that the carrier can be
distinguished from other carriers in the same system. Two distinct
counts must have distinct symbols. If two counts shared a symbol, the
ledger would conflate them. The adequacy condition formalizes this:
the carrier system must admit a decision procedure that, given two
symbols, decides whether they denote the same value.

\begin{leanexcerpt}{CarrierProcess}
structure CarrierProcess (Carrier: Type) where
  symbol : Fact
  value  : Number
  event  : Number $\to$ Number := ...
\end{leanexcerpt}

A \texttt{CarrierProcess} is parameterized by a type \texttt{Carrier} and
carries a \texttt{Fact} (the symbol), a \texttt{Number} (the value), and
an \texttt{event} transition rule. The structure is the chapter's minimum
admissible unit. Without the symbol, the value has no name in the ledger.
Without the event rule, the carrier cannot advance from one mark to the
next.

\begin{leanexcerpt}{DISTINGUISHABLE}
class DISTINGUISHABLE
    (Value : Type)
    (Observation : CarrierProcess Value) where
  fact : Fact
  symbol : Type Value
  distinct? : Prop := True
  different? : Type Value $\to$ Prop := ...
  dec_distinct : DecidablePred different?
\end{leanexcerpt}

A type \texttt{Value} together with a \texttt{CarrierProcess Value}
admits the \texttt{DISTINGUISHABLE} typeclass when the carrier supplies
a fact, a symbol at the appropriate universe level, an admissible
inequality predicate, and a decision procedure for that predicate. The
decision must terminate. Without \texttt{DISTINGUISHABLE}, the ledger
cannot reliably detect duplicate entries; with it, the ledger can.

For the integers, the decision procedure is trivial: compare the digits.
For polynomials, the procedure is more involved but still finite: subtract
the two polynomials and check whether all coefficients are zero. For
arbitrary functions, the procedure does not exist in general; two
functions agreeing on every input may differ in their internal
representation, and no algorithm can decide function equality from finite
data. The compiler is honest about this: it refuses to provide
\texttt{DISTINGUISHABLE} instances for function types without further
information.

The adequacy of a carrier is therefore not automatic. It must be proved
for each specific type. For most concrete types in mathematics
(naturals, rationals, polynomials, finite sets), the proof is straightforward.
For abstract types (function spaces, equivalence classes without canonical
representatives), the proof may fail. When it fails, the type cannot serve
as a carrier in this system.

The chapter's claim is that the carrier/value distinction, together with
the adequacy condition, is what makes the ledger trustworthy. A mark in
the ledger is a \texttt{CarrierProcess}; two marks can be compared by
the \texttt{DISTINGUISHABLE} decision procedure; equal marks are
identified; unequal marks remain distinct.

\section{Partial Order and Distinguishability}
\label{se:partial-order-and-distinguishability}

The divisibility relation on the positive integers is the chapter's
working example of a partial order that is not a total order. Two
positive integers $m, n$ satisfy $m \mid n$ when there exists a positive
integer $k$ such that $n = mk$. The relation is reflexive ($n \mid n$
because $n = n \cdot 1$), antisymmetric (if $m \mid n$ and $n \mid m$,
then $m = n$, because $m \le n$ and $n \le m$), and transitive (if
$m \mid n$ and $n \mid \ell$, then $m \mid \ell$).

The divisibility relation is not total. Six does not divide ten and ten
does not divide six; the pair $\{6, 10\}$ is incomparable. The pair has
a greatest common divisor (the meet): $\gcd(6, 10) = 2$. The pair has a
least common multiple (the join): $\operatorname{lcm}(6, 10) = 30$. Both
exist; neither equals six nor ten.

The lattice of positive integers under divisibility has a least element
(the integer 1, which divides every other positive integer) and no
greatest element (the lattice is unbounded above). Every finite set of
positive integers has a meet and a join. The lattice is distributive:
$\gcd(a, \operatorname{lcm}(b, c)) = \operatorname{lcm}(\gcd(a, b),
\gcd(a, c))$.

The divisibility lattice differs structurally from the usual numerical
ordering. Under the usual order, $5 < 6 < 10 < 12$ is a chain. Under
divisibility, $5$ and $6$ are incomparable; $5$ and $10$ are comparable
($5 \mid 10$); $5$ and $12$ are incomparable; $6$ and $10$ are
incomparable; $6$ and $12$ are comparable ($6 \mid 12$); $10$ and $12$
are incomparable. The Hasse diagram of the divisibility relation is a
two-dimensional structure: prime factorization expressed graphically.

This is the chapter's instance of a non-total partial order. Each
element has a position in the order; the position is its prime
factorization; the relation is whether one factorization divides
another. The order records structure that a total order would
collapse: which factors are present, which are absent, how they
combine.

In Lean, the divisibility relation is the predicate \texttt{Nat.dvd}
in Mathlib. The type \texttt{Number} in \texttt{Episode1.lean} can
carry either the standard order or divisibility as its
\texttt{Number.le} instance, depending on the typeclass resolution.
Both are valid partial orders; the typeclass cascade picks one
according to context.

The DISTINGUISHABLE typeclass plays a different role here. Two integers
that compare incomparable under divisibility are still distinguishable as
integers: $6$ and $10$ are different values even if neither divides the
other. The DISTINGUISHABLE relation is a finer relation than divisibility.
It is the underlying equality of the carrier. The divisibility relation
is a coarser structure built on top.

This is the chapter's third claim. The partial order is built from
DISTINGUISHABLE: two elements are comparable under the order only
because they are first distinguishable as elements. The order is
secondary; the distinguishability is primary. Different orders can be
built on the same set of distinguishable elements; each captures a
different aspect of the elements' relationships.

\section{Cauchy Sequences and the Limit}
\label{se:cauchy-sequences-and-the-limit}

The decimal expansion of $\pi$ begins $3, 3.1, 3.14, 3.141, 3.1415,
3.14159, \ldots$. Each term is a rational number. The sequence has no
terminating last entry; there is always a next decimal place. Yet the
sequence is in a precise sense convergent: the difference between the
$n$-th term and the $m$-th term shrinks toward zero as both $n$ and $m$
grow large.

This is the Cauchy condition. A sequence $(a_n)$ of rationals is
\emph{Cauchy} when for every positive rational $\varepsilon$, there
exists $N$ such that $|a_n - a_m| < \varepsilon$ for all $n, m \ge N$.
The condition is internal to the rationals: it does not assume the
existence of a limit. It only says the terms are clustering.

The expansion of $\pi$ is Cauchy: the difference between the $n$-th and
$m$-th terms is at most $10^{-\min(n,m)}$, which shrinks geometrically.
Yet there is no rational number to which the sequence converges. Every
rational $q$ has the property that $|\pi - q| > 0$; the sequence
approaches no rational target. The Cauchy sequence has a gap to fill.

The gap is filled by the construction of the real numbers. The reals
are defined as equivalence classes of Cauchy sequences of rationals,
where two sequences are equivalent if their difference converges to
zero. The real number $\pi$ is the equivalence class of the decimal
expansion (and of any other rational sequence that converges to the
same limit). The reals are the completion of the rationals under the
Cauchy condition.

This is the chapter's fourth claim: the partial order on rationals,
together with the Cauchy condition, forces the existence of the reals.
The reals are not assumed; they are derived. A ledger of rational marks
that admits Cauchy sequences as limits must extend to include real
numbers as the limits' values. The extension is not optional.

The Lean construction is precisely this. The type \texttt{Real} in
\texttt{Mathlib} is defined as the quotient of Cauchy sequences of
rationals by the equivalence relation of approaching the same limit.
The \texttt{Limit} type in \texttt{Episode1.lean} carries this:
a \texttt{Limit} is a sequence together with a rational target. The
\texttt{LimitProcess} structure carries the indexing and the iterate
rule; the \texttt{ENCODED} certificate marks the limit as carried.

\begin{leanexcerpt}{LimitProcess}
structure LimitProcess
    (Value : Type)
    (Carrier : CarrierProcess Value)
    [DISTINGUISHABLE Value Carrier]
    [ADMISSIBLE Value Carrier]
    [count : COUNTABLE Value Carrier]
  where
  indexing_process : IndexingProcess Value Carrier
  limit    : Rational
  sequence : Sequence
  iterate  : Sequence $\to$ Sequence := ...
\end{leanexcerpt}

The \texttt{LimitProcess} typeclass requires two pieces of data: a function
that produces the $n$-th approximation, and a function that gives, for
each positive $\varepsilon$, an index $N$ beyond which all approximations
agree to within $\varepsilon$. Together they specify a Cauchy sequence
without yet specifying a limit.

\texttt{ENCODED} certifies that the \texttt{LimitProcess} has reached a
state where the limit is admissible. The limit may not be a rational; it
may be a value that the rational record cannot contain directly. The
encoding is the assertion that the limit is named, even if its precise
value is not yet computable.

The Cauchy completion of the rationals is the mathematical structure
that the chapter's ledger requires. Rational marks are not sufficient
to carry all admissible records, because admissible sequences may have
irrational limits. The completion extends the ledger automatically. The
extension is forced by the Cauchy condition; it is not a choice.

\section{The Residue}
\label{se:the-residue}

Divide the polynomial $x^3 - 2$ by the linear polynomial $x - 1$. The
quotient is $x^2 + x + 1$ with remainder $-1$:
\[
x^3 - 2 = (x - 1)(x^2 + x + 1) + (-1).
\]
The remainder is the residue: what is left over after the admissible
division is performed. The residue is not an error. It is a structured
leftover. The remainder theorem says: the residue of dividing $p(x)$ by
$x - a$ is exactly $p(a)$. In this case, $p(1) = 1 - 2 = -1$.

The residue carries the part of $p(x)$ that the divisor $x - 1$ cannot
reach. The quotient $x^2 + x + 1$ is the part that $x - 1$ reaches.
Together they reconstruct $p(x)$. Each carries a different aspect of
the original polynomial.

The residue concept generalizes far beyond polynomials. In any
algebraic structure with a notion of division, the residue is what
remains. In number theory, dividing 17 by 5 gives quotient 3 and
residue 2. In ring theory, the residue ring $\mathbb{Z}/n\mathbb{Z}$
is the structure of residues modulo $n$. In analysis, the residue
theorem of complex analysis names the value left at a pole of a
meromorphic function.

In each case, the residue is the part that the division does not
discharge. The admissible quotient is the structured part; the
residue is what remains. Neither alone reconstructs the original; both
together do.

The chapter's fifth claim: the residue is part of the record. A
ledger that records a division operation must record both the
quotient and the residue. Discarding the residue would lose
information; reporting only the quotient would falsely claim
that the division is exact when it is not.

In Lean, the residue appears as the \texttt{RESIDUE} typeclass in
\texttt{Episode1.lean}:

\begin{leanexcerpt}{RESIDUE}
class RESIDUE
    (Value : Type)
    (Carrier : CarrierProcess Value)
    [d : DISTINGUISHABLE Value Carrier]
    [a : ADMISSIBLE Value Carrier]
    [c : COUNTABLE Value Carrier]
    [e : ENCODED Value Carrier]
  where
  cauchy_process : CauchyProcess Value Carrier
  representative? : Limit $\to$ Limit $\to$ Prop := ...
\end{leanexcerpt}

A \texttt{RESIDUE} instance carries a Cauchy process and a
predicate that says when one limit-state is the same residue as
another. The class is admissible only after the prior cascade
(\texttt{DISTINGUISHABLE}, \texttt{ADMISSIBLE}, \texttt{COUNTABLE},
\texttt{ENCODED}) has been provided. The compiler enforces this
stack at elaboration time: a missing instance below blocks
construction above.

The \texttt{Sample} typeclass extends the residue concept to
measurement. A \texttt{Sample} is a measured value together with
its residue: the part the instrument captured and the part it
missed. The \texttt{Sample} is the algebraic shadow of the
physical-measurement notion. The chapter's mathematical content
is independent of any physical measurement; the algebraic structure
holds whether the values come from polynomials, integers, or
measurements.

The chapter closes here. The smallest admissible mark a ledger can
carry is a triple: a value (the fact), a carrier (the symbol that
names it), and a partial order (the relation among carriers). The
adequacy condition is DISTINGUISHABLE: the carrier must admit a
decision procedure for equality. The order is built on top of the
adequacy condition. Cauchy completion extends the ledger
automatically when sequences require limits. The residue carries
what division does not discharge.

These five structures --- set, carrier, partial order, Cauchy
completion, residue --- together constitute the chapter's answer
to its closed question.

\begin{bridge}

The chapter's question was what the smallest admissible mark a ledger can
carry is.

The answer is a triple: a value, a carrier symbol that names it, and a
partial order that locates it among other values. The carrier must
admit a decision procedure for equality (\texttt{DISTINGUISHABLE}). The
partial order may be built on top of the decision procedure (the
divisibility lattice is one example). Cauchy sequences in the ordered
record may have limits outside the original set; the completion is
automatic. The residue carries what division cannot discharge.

What remains is comparison. The chapter has named the marks and their
order, but has not yet examined what it means for an instrument to
compare two marks and produce a repeatable result. The next chapter
asks how comparison becomes admissible: not how marks differ, but how
an instrument can carry the difference reliably without confusing the
comparison with the marks themselves.

\end{bridge}

\begin{coda}{The Archive Accession Desk}

The accession desk is the first station any incoming document
crosses on its way into the archive. A new collection arrives this
morning: forty letters from a deceased faculty member, two notebooks,
a folder of typed manuscripts, and an unopened envelope marked
``personal correspondence, do not file without consultation.''
Conservation has cleared it; the desk now processes it.

The archivist's first task is to give each item a unique accession number.
The accession number is the mark. Without it, the item exists in the world
but has no entry in the archive's catalog: it is unfindable. With the
accession number, the item is admissible to the archive's records.

The accession number for the first letter is M-2026-1847-001. The format is
specified by the archive's policy: M for manuscript collection, 2026 for the
year of accession, 1847 for the collection's sequence number within the year,
001 for the first item in the collection. The number is unique: no two items
in the archive share an accession number, and no two collections share a
sequence number within a given year.

This is the carrier in the chapter's sense. The letter is the value; the
accession number is the symbol that names it. The number distinguishes this
letter from every other letter in the archive. Reading the number alone
tells you nothing about the letter's content: not the sender, not the
recipient, not the date, not the topic. The number is the locator; the
content is what it locates.

The number must be \texttt{DISTINGUISHABLE}. The format guarantees this: two
different items receive two different numbers. The archive's database
enforces uniqueness: an attempt to register two items with the same number
produces an error. The decision procedure for equality is trivial: compare
the numbers digit by digit.

The archivist enters the first letter into the database. The entry has
several fields: the accession number, the brief description (sender,
recipient, date, topic), the location code (where in the archive the letter
is physically stored), the access restriction code (open, restricted,
sealed), and the date of accession.

The brief description is the catalog's surface: what a researcher sees when
they search the database. The location code tells the staff where to
retrieve the letter physically. The access code tells them whether the
researcher may consult it. The accession date is the timestamp of when the
letter became admissible to the archive.

The letter's relation to other letters in the archive is the partial order.
Letters from the same correspondent are related by author. Letters from the
same date are related by chronology. Letters about the same topic are
related by subject. None of these relations is total: many letters have
different correspondents, different dates, different topics; not every pair
of letters is comparable along every axis. The archive's catalog admits
multiple partial orders simultaneously: search by author, search by date,
search by topic. Each is a different lens on the same set of items.

This is the chapter's partial-order structure. The archive does not have a
single global ranking of letters; it has a collection of admissible
relations. A researcher's query selects one. The result is the letters
that compare favorably under that relation.

After the first letter, the archivist proceeds through the rest of the
collection. The forty letters are entered one by one. The two notebooks
receive accession numbers 041 and 042. The typed manuscripts receive
numbers 043 through 067 (twenty-five manuscripts, each a separate
accession). The sealed envelope receives number 068; it is flagged with the
access code ``RESTRICTED --- consult curator''. The collection's total is
sixty-eight items.

The Cauchy completion concept appears in a quieter form. The archive has a
policy of allowing supplementary materials to be added to a collection
after initial accession. If a researcher later discovers a related
document --- a letter that completes a chain, a manuscript that includes a
missing draft --- the document can be added with a sequence number
continuing where the original accession left off. The collection is a
Cauchy sequence: each new item refines the collection's representation,
without disturbing the prior items.

In principle, the collection's complete representation could require
infinitely many supplementary additions: every reference in every letter
might lead to a further document that should be included. In practice, the
archive closes the collection at some point and considers it complete. The
closure is a Cauchy condition: the supplementary additions become smaller
and less essential, until the marginal value of additional items is below
the archive's threshold. The collection is then sealed at its current
sequence number.

The residue concept appears too. Some items in the collection cannot be
processed into the standard accession format. The unopened personal
correspondence requires additional consultation; it is held in a temporary
restricted folder while the donor's family decides. The deceased
colleague's notes in the margins of the typed manuscripts include initials
that have not been identified; those notes carry a metadata residue (the
unidentified initials) that future research may resolve. The collection
admits these residues without erasing them; they are noted as
\texttt{UNRESOLVED} in the database.

By late afternoon, the collection has been processed. Sixty-seven items
have been entered into the database with full metadata. One item (the
sealed envelope) has been flagged for curator consultation. The collection's
overall description has been drafted by the archivist and submitted to the
collection's named curator for review. The collection has been assigned to
shelf positions in the third-level stacks.

A graduate student stops by the desk in the late afternoon. She is
researching a topic adjacent to the deceased colleague's work and has heard
that a new collection has been accessioned. She asks whether she may
consult the collection. The archivist explains that the collection is
still in preliminary processing; the brief descriptions have been entered,
but the items are not yet fully cataloged. The student may search the brief
descriptions but cannot yet request specific items by their accession
numbers; the items are not yet physically retrievable.

This is the difference between accession and full processing. The accession
mark is admissible immediately: the items are now part of the archive's
records. Full processing --- detailed descriptions, subject indexing,
finding aids --- takes additional weeks. The accession mark is the minimum;
full processing is the comfortable middle. Neither is a final state; the
collection's record will continue to be enriched as researchers consult it,
as related documents are discovered, as the curator updates the metadata.

The accession desk's discipline matches the chapter's claim. Every item
that enters the archive receives a mark. The mark is the smallest unit the
ledger can carry. The mark is \texttt{DISTINGUISHABLE} from all other marks.
The collection's items are related by multiple partial orders (author, date,
topic). Supplementary additions extend the collection as a Cauchy
refinement. The residues (unresolved initials, sealed envelopes) are
carried explicitly as unresolved entries.

At five o'clock the archivist closes the desk for the day. The collection
is now part of the archive. Tomorrow she will process another collection
that has been waiting in the incoming queue. The discipline of the
accession desk is the chapter's mathematical structure in institutional
form: every item is a mark, every mark is admissible, every mark is part
of the record.

\end{coda}
